Una persona de masa $100\,\mathrm{kg}$ se deja caer desde $3000\,\mathrm m$, con velocidad vertical inicial nula. La resistencia del aire es proporcional a la velocidad, con coeficiente $20\,\mathrm{kg/s}$ antes de abrir el paracaídas y $100\,\mathrm{kg/s}$ después. Use $g=9.8\,\mathrm{m/s^2}$, suponga apertura instantánea sin impulso y mida los tiempos desde el inicio de la caída.

\begin{enumerate}[label=(\alph*),ref=\alph*,leftmargin=*]
\item ¿Cuándo alcanza el suelo si abre el paracaídas a los $30$ segundos?
\item ¿Cuándo alcanza el suelo si la apertura se retrasa hasta los $60$ segundos?
\end{enumerate}
